\section{What is proved and what is not}\label{sec:end}

\subsection*{Proved}
The following hold without any unproved hypothesis.
\begin{enumerate}[label=\textup{(\arabic*)}]
\item The Hodge conjecture for the Fermat variety $X^{n}_{m}$ of every
dimension and every degree $m=2^{a}3^{b}19^{c}$ with $a\le1$, for every
product of Fermat varieties of such a degree, and for every abelian variety
of Fermat type of such a degree (\Cref{thm:fermat114}). Besides the results
of this paper, the proof uses Aoki's reduction \cite{Aok83,Aok00}, Shioda's
inductive structure \cite{Shi79}, Aoki's cycles \cite{Aok87} and the
Lefschetz theorem on $(1,1)$-classes, and it rests on one certified
computation in interval arithmetic (\Cref{prop:family2} and \Cref{app:cert114}). It is
new when $57\mid m$ and $m\ne57$; for $m=57$ the earlier proof \cite{MMRV26}
depends on the statement discussed in \Cref{rem:kang}.
\item The Hodge conjecture for the Fermat fourfold $X^{4}_{m}$ of every odd
degree $m$ (\Cref{thm:oddfour}); for $m$ prime to $6$, the classification of the
Hodge characters behind it: each contains a pair or is $5$-standard
(\Cref{thm:sextuples}); and for $3\mid m$, the reduction of every Hodge
character without a pair, by moves, to three exceptional orbits
(\Cref{thm:moves} and \Cref{lem:exceptional}). Besides the results of this paper, the proof uses
the inductive structure of Shioda and Ran, Aoki's cycles \cite{Aok87}, and
the classical nonvanishing of $B_{1,\chi}$ for odd primitive Dirichlet
characters \cite{Was97}.
\item The equivalences of \Cref{thm:equiv}, and the description in \Cref{thm:routes}
and \Cref{cor:bypass} of every set of the statements of \Cref{tab:rules}
from which the conjecture follows by those implications.
\item The statement \st{F3$'$} for every Fermat variety, every product of
Fermat varieties and every variety onto which such a product maps
surjectively (\Cref{thm:fermatcm}), together with the description of the Hodge classes of Fermat varieties by
characters (\Cref{prop:fermatchars}).
\item The closedness of the bounded-degree loci $\Sigma_{d}$ in a split Weil
family, and the criterion of \Cref{thm:weilloci} for $\Sigma=S$.
\item For every \'etale cyclic triple cover of a curve of genus $n+1\ge3$:
its Prym variety is of split Weil type (\Cref{lem:prym}); Schoen's
construction gives an irreducible $n$-dimensional subvariety $Y$ whose class
pairs nontrivially with the Weil plane; at a member with maximal Hodge group
$\cl(Y)=c\,\eta^{n}+w$ with $c>0$ and $w\ne0$ (\Cref{thm:schoenY}); and the Prym
variety is Weil-algebraic (\Cref{cor:schoen}, first proved by Schoen).
\item The transfer of \Cref{thm:transfer}: an abelian variety of Weil type for
$\QQ(\sqrt{-3})$ is, up to isogeny, a factor $A\times A'$ of a Prym variety,
and $A$ is Weil-algebraic if and only if $A'$ is.
\item The count of \Cref{prop:branched}: Prym varieties of cyclic triple
covers, \'etale or branched, form families of dimension at most $3n$, so for
$n\ge4$ they miss the very general member of the split Weil family.
\item The rigidity of \Cref{prop:abelprym}: the Abel--Prym curve of an
\'etale cyclic triple cover deforms with its Prym variety only along the
Prym locus, although its class stays algebraic on the whole split Weil
family; so for $n\ge4$ and maximal Hodge group, when it is embedded, it is
not semiregular.
\item The description of $Y$ in \Cref{prop:prymbn}: up to an \'etale map,
it is a component of the locus of line bundles with a nonzero section in
the fibre of the norm map over $[K_{X}]$. It is smooth of dimension $n$ at
its general point, with an explicit tangent space.
\end{enumerate}

\subsection*{Proved under a stated hypothesis}
\begin{enumerate}[label=\textup{(\arabic*)},resume]
\item Under \st{CM}: the Hodge conjecture for every Fermat variety, every
product of Fermat varieties and every variety such a product dominates
(\Cref{thm:fermatcm}), and the algebraicity of the Weil classes at every CM point
of a split Weil family (\Cref{thm:weilloci}).
\item Under the semiregularity of one subvariety $Y$ of genus $n+1$ with
maximal Hodge group: the Weil classes of every abelian variety of split Weil
type for $\QQ(\sqrt{-3})$ of dimension $2n$ (\Cref{cor:schoensemireg}).
\end{enumerate}

\subsection*{Not proved}
The Hodge conjecture; each of \st{F2}, \st{F3$'$}, \st{L}, \st{M}, \st{V},
\st{IP}; the semiregularity of $Y$ for any $n\ge4$
(\Cref{q:schoen}); \st{CM}, which is claimed in the unrefereed preprint
\cite{OAI26} and used here only as a hypothesis; and the Hodge conjecture for
the Fermat fourfolds of degrees $110$ and $220$, the only degrees up to $250$
left open by \Cref{thm:fermat114,thm:oddfour}, \cite{Pet26} and the census
\cite{Jum26b}.

\subsection*{The first open cases}
By \Cref{thm:routes} every route passes through \st{F2} and \st{F3$'$}, and each
has a first case that no known method reaches.

\emph{Weil classes in dimensions six and eight.} \st{F2} holds in dimension
at most five by \cite{MZ99} and the Weil fourfolds of the preprint
\cite{Mar25a} (\Cref{ssec:weilknown}). In dimension six the Weil classes of
split forms are algebraic, for $\QQ(\sqrt{-3})$ by \cite{Sch88,Sch98} and
in general by \cite{Mar25a}, and those of nonsplit forms are open.
For $2n=8$ we know of no family of Weil type whose general member has
maximal Hodge group and which carries algebraic Weil classes at every
member. For the split family of
$\QQ(\sqrt{-3})$ this is \Cref{q:schoen} with $n=4$, where $Y$ would have to
deform in $n^{2}-3n=4$ directions beyond the Prym locus.
\Cref{prop:branched} shows that Prym varieties of branched cyclic triple
covers do not reach the general member either, and \Cref{rem:count} shows
the same for \Cref{cor:transfer}. The Abel--Prym curve, built from the same
map as $Y$, does not deform in those directions (\Cref{prop:abelprym}), while
the Prym theta divisor of an \'etale double cover, the analogue of $Y$ by
\Cref{prop:prymbn}, deforms with every principally polarized abelian
variety; neither decides the question for $Y$. Preprints announcing the split eightfolds are cited in
\cite{BhH8}; we could not locate them in the public release of
\cite{OAI26}, and we do not use them.

\emph{Beyond Weil classes.} On the square $A\times A$ of a Mumford fourfold
$A$ \cite{Mum69} there are Hodge classes that are not generated by divisor
classes and Weil classes. They are the first instance of \st{F2} that the
Weil route does not cover. In \cite[Theorem~21.4]{BhH8} they are reduced to
the Kuga--Satake class of one K3 surface, which is not known to be
algebraic.

\emph{The first case of \st{F3$'$}.} For a K3 surface $S$ whose transcendental
lattice has a totally real field of degree at least two as its algebra of
Hodge endomorphisms \cite{vG08}, the classes on $S\times S$ that induce these
endomorphisms are not known to be algebraic in general \cite{BvGS24}; $S\times
S$ lies in the reach of abelian varieties by algebraic correspondences once
the Kuga--Satake correspondence of $S$ is algebraic
\cite[Remark~20.36]{BhH8}. They are algebraic on several explicit families,
where the real multiplication comes from double planes or from isogenies
\cite{Sch10,vGS25}, and on the general member of the first
four-dimensional families with real multiplication by a real quadratic
field, through symplectic automorphisms \cite{Var23}. For real multiplication by $\QQ(\sqrt d)$ it would
suffice to find an algebraic correspondence from $S$ to some K3 surface $S'$
that induces a Hodge similarity $T(S)\to T(S')$ of multiplier $d$: composing
its inverse with $\sqrt d$ gives a rational Hodge isometry, which is
algebraic by the theorem of Buskin and Huybrechts \cite{Bus19,Huy19}. For a real quadratic field the
maximal families have Picard number $2$ and dimension $8$ \cite{vG08}, and
their very general member is the first open case.

So the conjecture is equivalent to two statements, and each has an explicit
first open case. What this paper adds to the Weil case is that, for
$\QQ(\sqrt{-3})$ in every dimension, it now depends on a single, explicitly
constructed subvariety of a Prym variety and the question whether it is
semiregular.

\subsection*{A note on preparation}
This paper was prepared with substantial help from an AI assistant, working
under the author's direction; the author checked the arguments and is
responsible for the content. The assistant took part in every section, and
we list the main ideas it proposed. In \Cref{sec:fermat114}: the extension
of Peterson's method to block families (\Cref{ssec:families}), the
characters $\alpha_{1},\alpha_{2}$ and the join argument
(\Cref{ssec:twochars,ssec:join}), and the two families with the certificate
of \Cref{app:cert114}, found by computer searches that the assistant
designed. In \Cref{sec:fermat}: the classification of Hodge sextuples
(\Cref{thm:sextuples}) and the moves of \Cref{thm:moves}. In
\Cref{sec:weil}: the twisted-cohomology computation behind
\Cref{thm:schoenY} and the description of $Y$ in \Cref{prop:prymbn}. In
\Cref{sec:routes}: the enumeration of routes in \Cref{thm:routes}. The
assistant also drafted much of the text and wrote the programs that check
the finite steps.
